\documentclass[11pt,a4paper]{article}

\usepackage[provide=*,spanish]{babel}
\usepackage[utf8]{inputenc}
\usepackage[T1]{fontenc}
\usepackage{lmodern}
\usepackage[margin=2.35cm]{geometry}
\usepackage{amsmath}
\usepackage{booktabs}
\usepackage{longtable}
\usepackage{tabularx}
\usepackage{array}
\usepackage{enumitem}
\usepackage{listings}
\usepackage{xcolor}
\usepackage{microtype}
\usepackage{parskip}
\usepackage{titlesec}
\usepackage{hyperref}
\usepackage{caption}
\usepackage{float}

\definecolor{azulWCD}{RGB}{17,67,112}
\definecolor{azulClaro}{RGB}{232,241,249}
\definecolor{grisCodigo}{RGB}{247,247,247}
\definecolor{verdeCodigo}{RGB}{25,110,55}
\definecolor{rojoCodigo}{RGB}{145,35,35}
\definecolor{tituloColor}{RGB}{0,55,110}
\definecolor{reglaColor}{RGB}{0,90,160}

\hypersetup{
  colorlinks=true,
  linkcolor=azulWCD,
  urlcolor=azulWCD,
  citecolor=azulWCD,
  pdftitle={Estructura de los archivos de simulacion para el analisis de moderacion y captura de neutrones},
  pdfauthor={Jaime A. Betancourt}
}

\lstdefinestyle{tsvstyle}{
  backgroundcolor=\color{grisCodigo},
  basicstyle=\ttfamily\scriptsize,
  breaklines=true,
  breakatwhitespace=false,
  frame=single,
  rulecolor=\color{gray!40},
  tabsize=2,
  showstringspaces=false,
  columns=fullflexible,
  keepspaces=true
}
\lstset{style=tsvstyle}

\titleformat{\section}{\normalfont\Large\bfseries\color{azulWCD}}
  {\thesection.}{0.6em}{}
\titleformat{\subsection}{\normalfont\large\bfseries}
  {\thesubsection.}{0.5em}{}
\setlist{itemsep=0.25em,topsep=0.4em}
\renewcommand{\arraystretch}{1.18}
\setlength{\emergencystretch}{2em}

\newcommand{\codigo}[1]{\nolinkurl{#1}}
\newcommand{\archivo}[1]{\texttt{#1}} % \texttt en vez de \path: \path es frágil dentro de \caption{}
\newcommand{\var}[1]{\texttt{#1}}

\begin{document}
\pagestyle{empty}

%------------------------------------------------
% PORTADA
%------------------------------------------------

\begin{titlepage}
\begin{center}

\vspace*{0.5cm}
\textcolor{reglaColor}{\rule{\linewidth}{1.2pt}}

\vspace{1.3cm}
{\Large\scshape Universidad Industrial de Santander}\\[0.2cm]
{\large Escuela de Física --- Facultad de Ciencias}\\[0.2cm]
{\large Doctorado en Física}

\vspace{2.2cm}
{\Huge\bfseries\color{tituloColor} Estructura de los archivos\\[0.2cm]
de simulación del WCD}

\vspace{1.8cm}
{\large\textbf{Informe técnico}}

\vspace{0.6cm}
\begin{minipage}{0.85\linewidth}
\centering
\textit{``Diccionario de datos para el análisis de moderación, captura, reflexión
y transmisión de neutrones''}
\end{minipage}

\vfill

{\large\textbf{Autor}}\\[0.15cm]
Jaime A. Betancourt\\[0.1cm]
Doctorado en Física, Universidad Industrial de Santander

\vspace{1.2cm}
\textcolor{reglaColor}{\rule{0.5\linewidth}{0.6pt}}

\vspace{1cm}
15 de septiembre de 2026

\vspace{1cm}
\textcolor{reglaColor}{\rule{\linewidth}{1.2pt}}

\end{center}
\end{titlepage}

\pagestyle{plain}

\begin{abstract}
Este informe documenta, columna por columna, la estructura de los cuatro archivos TSV que
produce la simulación Geant4/MEIGA del detector Cherenkov de agua (WCD) y que alimentan el
análisis de moderación y captura de neutrones descrito en
\textit{Informe-tecnico-Analisis-Moderacion-Captura-Neutrones}: \archivo{neutrones-incidentes.tsv},
\archivo{pasos-neutrones.tsv}, \archivo{interacciones-neutrones.tsv} y
\archivo{capturas-neutrones.tsv}. Es un documento de referencia --un diccionario de datos-- para
cualquier persona que necesite leer, depurar o extender ese análisis sin tener que releer el
código o volver a inspeccionar los archivos desde cero. Se describen las convenciones comunes a
los cuatro archivos, el significado físico y las unidades de cada columna, ejemplos reales de
filas, los valores categóricos que puede tomar cada campo, y cómo se relacionan entre sí mediante
el identificador \var{run\_id}.
\end{abstract}

\tableofcontents

\section{Objeto y alcance del informe}

Cada corrida de simulación (agua pura o agua con una concentración dada de NaCl) produce, entre
otros, cuatro archivos de texto separados por tabuladores (TSV) que registran, paso a paso, el
transporte de cada neutrón primario a través del detector. Estos cuatro archivos son la
\emph{única} fuente de datos del análisis de moderación y captura documentado en el informe
técnico complementario; este documento describe su estructura interna --nombres de columna, tipos
de dato, unidades y significado físico-- para que puedan leerse, verificarse o reutilizarse sin
depender de un notebook o una conversación previa. El análisis exploratorio que originó este
informe vive en el \textit{notebook} \codigo{Estructura-archivos-simulacion.ipynb}, que inspecciona
cada archivo con \codigo{pandas} (dimensiones, tipos, valores únicos, estadística descriptiva).

\section{Convenciones comunes a los cuatro archivos}
\label{sec:convenciones}

\begin{itemize}
  \item \textbf{Formato}: texto plano, columnas separadas por tabulador (\texttt{\textbackslash t}).
  La primera línea es una cabecera que empieza con \texttt{\#} seguido de los nombres de columna
  (también separados por tabulador). Al leer con \texttt{pandas.read\_csv(..., sep="\textbackslash t")}
  sin más argumentos, el símbolo \texttt{\#} queda pegado al nombre de la primera columna
  (\texttt{"\#run\_id"}) y debe limpiarse explícitamente (\codigo{columns.str.lstrip("\#")}) o, de
  forma más robusta, leerse con \codigo{skiprows=1} y una lista de nombres de columna ya definida.
  \item \textbf{Identificador de trazabilidad, \var{run\_id}}: entero de 0 a 99\,999 en una corrida
  de 100\,000 neutrones primarios. A pesar de su nombre, \textbf{no} corresponde al concepto de
  ``run'' de Geant4: es, en la práctica, un identificador único por neutrón primario inyectado, y es
  la columna que permite encadenar la trayectoria de un mismo neutrón entre los cuatro archivos.
  \item \textbf{Columnas degeneradas en estos archivos}: \var{event\_id} es constante (0) en las
  cuatro tablas; \var{track\_id} es constante (1, el neutrón primario) y \var{parent\_id} es
  constante (0, no tiene partícula padre) en los cuatro archivos, porque todos ellos ya vienen
  filtrados a la partícula \var{"neutron"} primaria. Estas columnas se conservan por compatibilidad
  con el formato genérico de salida de MEIGA, pero no aportan información variable en este contexto.
  \item \textbf{Unidades}: posiciones en centímetros (\var{\_cm}), energías cinéticas en MeV
  (\var{\_MeV}), tiempo global acumulado desde el inicio del evento en nanosegundos
  (\var{global\_time\_ns}).
  \item \textbf{Tamaño y volumen de datos} (corrida de referencia, agua pura, 100\,000 neutrones a
  25\,meV): ver la tabla~\ref{tab:tamanos}. Los dos archivos de historial completo
  (\archivo{pasos-neutrones.tsv} e \archivo{interacciones-neutrones.tsv}) superan los 500\,MB cada
  uno; se recomienda especificar \codigo{dtype} (p.\,ej. \codigo{category} para las columnas de
  texto) al cargarlos con \codigo{pandas} para reducir el uso de memoria.
\end{itemize}

\begin{table}[H]
\centering
\begin{tabularx}{\linewidth}{X r r}
\toprule
\textbf{Archivo} & \textbf{Filas (sin cabecera)} & \textbf{Tamaño en disco} \\
\midrule
\archivo{neutrones-incidentes.tsv} & 100\,000 & $\sim$15\,MB \\
\archivo{pasos-neutrones.tsv} & 3\,843\,486 & $\sim$668\,MB \\
\archivo{interacciones-neutrones.tsv} & 3\,243\,996 & $\sim$570\,MB \\
\archivo{capturas-neutrones.tsv} & 22\,748 & $\sim$3.6\,MB \\
\bottomrule
\end{tabularx}
\caption{Dimensiones de los cuatro archivos para la corrida de referencia (agua pura, 100\,000
neutrones, 25\,meV). El número de filas de \archivo{capturas-neutrones.tsv} varía con el medio
(23\,270 para agua con 2.5\,\% de NaCl, por ejemplo), porque depende de cuántos neutrones se
capturan en total.}
\label{tab:tamanos}
\end{table}

\section{\archivo{neutrones-incidentes.tsv}: registro de inyección}

Contiene exactamente \textbf{una fila por cada neutrón primario inyectado} (100\,000 filas): su
primer paso de transporte, siempre en material \var{Air}, con \var{step\_process =
Transportation}. Es el universo de referencia contra el que se calculan todos los porcentajes del
análisis (coeficiente de captura, reflexión, transmisión).

\begin{table}[H]
\centering
\small
\begin{tabularx}{\linewidth}{l l X}
\toprule
\textbf{Columna} & \textbf{Tipo} & \textbf{Significado} \\
\midrule
\var{run\_id} & entero & Identificador único del neutrón primario (0--99\,999). \\
\var{event\_id} & entero & Constante (0) en este archivo. \\
\var{particle} & texto & Constante (\var{"neutron"}). \\
\var{track\_id} & entero & Constante (1, partícula primaria). \\
\var{parent\_id} & entero & Constante (0, sin partícula padre). \\
\var{step\_number} & entero & Constante (1): es siempre el primer paso de la trayectoria. \\
\var{creator\_process} & texto & Constante (\var{"none"}: la partícula primaria no fue creada por
ningún proceso, fue inyectada directamente). \\
\var{step\_process} & texto & Constante (\var{"Transportation"}). \\
\var{material} & texto & Constante (\var{"Air"}): el neutrón se inyecta 2\,m por encima del tanque. \\
\var{pre\_x\_cm}, \var{pre\_y\_cm}, \var{pre\_z\_cm} & flotante & Posición inicial del neutrón (cm).
\var{pre\_z\_cm} = 200 para todos (altura de inyección). \\
\var{post\_x\_cm}, \var{post\_y\_cm}, \var{post\_z\_cm} & flotante & Posición al final de este
primer paso (cm); \var{post\_z\_cm} = 133.124 (superficie superior del tanque) para todos, ya que
el primer paso de transporte llega hasta la tapa. \\
\var{pre\_energy\_MeV}, \var{post\_energy\_MeV} & flotante & Energía cinética antes/después del
paso (MeV); ambas iguales entre sí (2.5$\times10^{-8}$\,MeV = 25\,meV) porque un paso de puro
transporte no cambia la energía. \\
\var{global\_time\_ns} & flotante & Tiempo global transcurrido desde el inicio del evento (ns). \\
\bottomrule
\end{tabularx}
\caption{Columnas de \archivo{neutrones-incidentes.tsv}.}
\end{table}

\noindent Ejemplo real (una fila, columnas seleccionadas):
\begin{lstlisting}
run_id=0  particle=neutron  pre_z_cm=200  post_z_cm=133.124  pre_energy_MeV=2.499996299e-08  global_time_ns=305793.6085
\end{lstlisting}

\section{\archivo{pasos-neutrones.tsv}: historial completo}

Contiene \textbf{todos} los pasos de transporte de \textbf{todos} los neutrones primarios (no solo
los capturados), incluyendo los pasos de puro cruce de frontera entre materiales
(\var{step\_process = Transportation}). Es el archivo más grande de los cuatro porque registra
cada paso, no solo las interacciones físicas. Comparte exactamente el mismo esquema de columnas que
\archivo{neutrones-incidentes.tsv} (tabla~\ref{tab:esquema-pasos}), salvo que aquí
\var{step\_number} crece (1, 2, 3, \ldots) a lo largo de la trayectoria de cada neutrón y
\var{step\_process}/\var{material} varían de un paso a otro.

\begin{table}[H]
\centering
\small
\begin{tabularx}{\linewidth}{l l X}
\toprule
\textbf{Columna} & \textbf{Tipo} & \textbf{Significado} \\
\midrule
\var{run\_id} & entero & Identificador del neutrón (igual convención que en
\archivo{neutrones-incidentes.tsv}). \\
\var{event\_id}, \var{track\_id}, \var{parent\_id} & entero & Constantes (0, 1, 0) --ver
sección~\ref{sec:convenciones}. \\
\var{step\_number} & entero & Índice secuencial del paso dentro de la trayectoria del neutrón,
empezando en 1. \\
\var{creator\_process} & texto & Constante (\var{"none"}: partícula primaria). \\
\var{step\_process} & texto & Proceso físico de Geant4 ocurrido en ese paso. Valores observados:
\var{Transportation} (cruce de frontera, sin interacción física), \var{hadElastic} (dispersión
elástica), \var{nCapture} (captura radiativa), \var{neutronInelastic} (dispersión inelástica,
infrecuente a estas energías). \\
\var{material} & texto & Material de Geant4 donde ocurrió el paso (\var{Air},
\var{G4\_STAINLESS-STEEL}, \var{Tyvek\_HDPE}, o el material del medio activo --p.\,ej.
\var{Water\_TS\_H\_of\_Water} para agua pura, \var{SaltyWater\_NaCl\_2.5pct} para 2.5\,\% de NaCl). \\
\var{pre\_x/y/z\_cm}, \var{post\_x/y/z\_cm} & flotante & Posición antes/después del paso (cm). \\
\var{pre\_energy\_MeV}, \var{post\_energy\_MeV} & flotante & Energía cinética antes/después del
paso (MeV). En un paso \var{Transportation} son iguales; en \var{hadElastic} difieren (la
diferencia es la base del cálculo de la letargía $\xi$); en \var{nCapture},
\var{post\_energy\_MeV} = 0 (el neutrón desaparece). \\
\var{global\_time\_ns} & flotante & Tiempo global acumulado (ns); crece monótonamente a lo largo de
la trayectoria de cada neutrón. \\
\bottomrule
\end{tabularx}
\caption{Columnas de \archivo{pasos-neutrones.tsv} (idéntico esquema en
\archivo{interacciones-neutrones.tsv}).}
\label{tab:esquema-pasos}
\end{table}

\noindent Ejemplo real: los tres primeros pasos del neutrón \var{run\_id=0} (entra en aire, se
dispersa elásticamente en el acero, y vuelve a transportarse):
\begin{lstlisting}
run_id=0  step_number=1  step_process=Transportation  material=Air              pre_energy_MeV=2.4999963e-08  post_energy_MeV=2.4999963e-08
run_id=0  step_number=2  step_process=hadElastic       material=G4_STAINLESS-STEEL  pre_energy_MeV=2.4999963e-08  post_energy_MeV=1.5518026e-08
run_id=0  step_number=3  step_process=Transportation  material=G4_STAINLESS-STEEL  pre_energy_MeV=1.5518026e-08  post_energy_MeV=1.5518026e-08
\end{lstlisting}

\section{\archivo{interacciones-neutrones.tsv}: solo interacciones físicas}

Mismo esquema de columnas que \archivo{pasos-neutrones.tsv} (tabla~\ref{tab:esquema-pasos}), pero
\textbf{excluye} las filas con \var{step\_process = Transportation}: solo quedan las filas donde
ocurrió una interacción física real (\var{hadElastic}, \var{nCapture} o \var{neutronInelastic}).
Es, en la práctica, un filtro de \archivo{pasos-neutrones.tsv} --se puede reconstruir a partir de
él-- pero se entrega como archivo aparte porque es la fuente más directa para contar colisiones,
calcular $\xi$ y $N$, y localizar el evento de captura de cada neutrón sin tener que descartar
primero los pasos de puro transporte. Para la corrida de referencia (agua pura), la distribución de
\var{step\_process} es: \var{hadElastic} 3\,219\,440 filas, \var{nCapture} 22\,748 filas,
\var{neutronInelastic} 1\,808 filas --el número de filas \var{nCapture} coincide exactamente con el
número de filas de \archivo{capturas-neutrones.tsv}, lo cual sirve como verificación cruzada de
consistencia entre ambos archivos.

\section{\archivo{capturas-neutrones.tsv}: solo eventos de captura}

Subconjunto de \archivo{interacciones-neutrones.tsv} filtrado a \var{step\_process = nCapture}
únicamente, en \textbf{cualquier} material (no solo el medio activo: incluye capturas en acero,
Tyvek y aire). Cada fila es el evento final de la trayectoria de un neutrón: por eso
\var{post\_energy\_MeV} vale siempre 0. A diferencia de los otros tres archivos, incorpora dos
columnas adicionales que identifican el núcleo capturador.

\begin{table}[H]
\centering
\small
\begin{tabularx}{\linewidth}{l l X}
\toprule
\textbf{Columna} & \textbf{Tipo} & \textbf{Significado} \\
\midrule
\multicolumn{3}{l}{\textit{(mismas columnas que la tabla~\ref{tab:esquema-pasos}, sin \var{creator\_process}, más:)}} \\
\var{target\_Z} & entero & Número atómico del núcleo que capturó al neutrón (p.\,ej. 1 = H, 8 = O,
11 = Na, 17 = Cl, 24/26/28 = Cr/Fe/Ni del acero, 7 = N del aire). \\
\var{target\_A} & entero & Número másico del isótopo capturador. \\
\bottomrule
\end{tabularx}
\caption{Columnas adicionales de \archivo{capturas-neutrones.tsv} respecto al esquema común.}
\end{table}

\noindent Ejemplo real (captura de un neutrón sobre hidrógeno del agua):
\begin{lstlisting}
run_id=1  step_process=nCapture  material=Water_TS_H_of_Water  target_Z=1  target_A=1  post_energy_MeV=0  global_time_ns=591226.533
\end{lstlisting}

\noindent Filtrar \var{material} que contiene la palabra \var{"Water"} \textbf{y} \var{target\_Z}
dentro del conjunto de núcleos esperado para el medio activo de esa corrida (H/O para agua pura;
H/O/Na/Cl si hay NaCl) es, precisamente, el primer paso del análisis de moderación y captura --ver
el informe técnico complementario.

\section{Cómo se relacionan los cuatro archivos}

La figura conceptual del flujo de datos es la siguiente:

\begin{enumerate}
  \item \archivo{neutrones-incidentes.tsv} fija el universo de 100\,000 \var{run\_id} y su estado
  de inyección (posición, energía y tiempo iniciales).
  \item \archivo{pasos-neutrones.tsv} contiene, para cada uno de esos \var{run\_id}, la trayectoria
  completa (todas las filas de un mismo \var{run\_id}, ordenadas por \var{step\_number}, describen
  un único neutrón desde que se inyecta hasta que se captura o abandona el mundo simulado).
  \item \archivo{interacciones-neutrones.tsv} es el mismo conjunto de filas que
  \archivo{pasos-neutrones.tsv}, filtrado a las que representan una interacción física real.
  \item \archivo{capturas-neutrones.tsv} es, a su vez, el subconjunto de
  \archivo{interacciones-neutrones.tsv} con \var{step\_process = nCapture} --la última fila de la
  trayectoria de cada neutrón que efectivamente se captura (no todos lo hacen: algunos se
  transmiten o se reflejan, ver el informe complementario).
\end{enumerate}

La clave \var{run\_id} es la única forma de unir estos cuatro archivos de forma correcta; nunca debe
asumirse que el orden de las filas por sí solo basta para relacionar un archivo con otro (aunque en
la práctica las filas de un mismo \var{run\_id} suelen aparecer contiguas y en el mismo orden
relativo entre archivos, esto no está garantizado y no debe usarse como criterio de unión).

\section{Reproducibilidad}

La exploración que documenta este informe --dimensiones, tipos de dato, valores únicos, filas de
ejemplo-- se realizó con el \textit{notebook} \codigo{Estructura-archivos-simulacion.ipynb}, ubicado
junto a los archivos TSV de cada corrida. El análisis derivado que consume estos cuatro archivos
--coeficiente de captura, reflexión, transmisión, $N$, $\xi$, $\Delta t$ y $\tau$-- está documentado
en el \textit{notebook} \codigo{analisis\_moderacion\_captura\_neutrones.ipynb} y en el informe
técnico \textit{Informe-tecnico-Analisis-Moderacion-Captura-Neutrones}, disponibles ambos en este
mismo repositorio (\archivo{analysis/notebooks/} y \archivo{docs/technical-reports/},
respectivamente).

\section{Conclusiones}

Los cuatro archivos TSV de cada corrida de simulación forman una jerarquía de subconjuntos
encadenados por \var{run\_id}: inyección $\supseteq$ historial completo $\supseteq$ interacciones
físicas $\supseteq$ eventos de captura. Conocer esta relación --y el significado exacto de cada
columna, resumido en las tablas de este informe-- es el prerrequisito para poder leer, verificar o
extender con confianza el análisis de moderación y captura de neutrones del WCD.

\end{document}
